% ============================================================
%  EEDT Version 5 — Zenodo Revised Draft (2026-03-15)
%  GPS v8 Hardware Validation + Theoretical Strengthening
%  Compiler: pdflatex (twice)
%  Figures: fig1_gps_smoking_gun.pdf
%           fig2_device_comparison.pdf
%           fig3_data_collapse.pdf
%           fig4_followup.pdf
% ============================================================
\documentclass[10pt,twocolumn]{article}

% ── Packages ─────────────────────────────────────────────────
\usepackage[T1]{fontenc}
\usepackage[utf8]{inputenc}


\usepackage{colortbl}
\usepackage[top=2.0cm,bottom=2.0cm,left=1.6cm,right=1.6cm,
            columnsep=0.6cm]{geometry}
\usepackage{amsmath,amssymb,amsthm}
\usepackage{mathtools}
\usepackage{graphicx}
\usepackage{booktabs}
\usepackage{array}
\usepackage{xcolor}
\usepackage{hyperref}

\usepackage{enumitem}
\usepackage{caption}
\usepackage{subcaption}
\usepackage{algorithm}
\usepackage{algpseudocode}
\usepackage{listings}
\usepackage{cite}
\usepackage{fancyhdr}
\usepackage{titlesec}
\usepackage{abstract}

% ── Colors ───────────────────────────────────────────────────
\definecolor{NatureRed}{HTML}{E64B35}
\definecolor{NatureBlue}{HTML}{4DBBD5}
\definecolor{NatureGreen}{HTML}{00A087}
\definecolor{NatureGold}{HTML}{F39B7F}
\definecolor{DarkGray}{HTML}{333333}

% ── Hyperref setup ───────────────────────────────────────────
\hypersetup{
  colorlinks=true,
  linkcolor=NatureBlue,
  citecolor=NatureGreen,
  urlcolor=NatureBlue,
  pdftitle={EEDT: Entanglement-Enhanced Dynamical Tracking (Version 5)},
  pdfauthor={Takeshi Okuda}
}

% ── Section formatting ───────────────────────────────────────
\titleformat{\section}{\normalfont\large\bfseries\color{DarkGray}}
  {\thesection.}{0.5em}{}
\titleformat{\subsection}{\normalfont\normalsize\bfseries\color{DarkGray}}
  {\thesubsection}{0.5em}{}
\titleformat{\subsubsection}{\normalfont\normalsize\itshape}
  {\thesubsubsection}{0.5em}{}
\titlespacing{\section}{0pt}{8pt}{4pt}
\titlespacing{\subsection}{0pt}{6pt}{3pt}

% ── Caption / listings ───────────────────────────────────────
\captionsetup{font=small, labelfont={bf,color=DarkGray}, skip=4pt}
\lstset{
  basicstyle=\ttfamily\footnotesize,
  backgroundcolor=\color{gray!8},
  frame=single, framerule=0.4pt, rulecolor=\color{gray!40},
  xleftmargin=6pt, xrightmargin=6pt,
  breaklines=true, breakatwhitespace=true,
  tabsize=2, showstringspaces=false, columns=flexible
}

% ── Theorem environments ─────────────────────────────────────
\newtheorem{result}{Analytical Result}
\newtheorem{proposition}{Proposition}
\newtheorem{corollary}{Corollary}
\newtheorem{remark}{Remark}
\newtheorem{definition}{Definition}

% ── Math shortcuts ───────────────────────────────────────────
\newcommand{\GPS}{G_{\mathrm{PS}}}
\newcommand{\nuZZ}{\nu_{\mathrm{ZZ}}}
\newcommand{\oZZ}{\omega_{\mathrm{ZZ}}}
\newcommand{\Tmc}{T_{\mathrm{MCM}}}
\newcommand{\taustar}{\tau^*}
\newcommand{\phistar}{\phi^*}
\newcommand{\Fstro}{\mathcal{F}_{\mathrm{stro}}}
\newcommand{\Fth}{\mathcal{F}_{\infty}}
\newcommand{\Fref}{\mathcal{F}_{\mathrm{ref}}}
\newcommand{\Nmeas}{N_{\mathrm{meas}}}
\newcommand{\eps}{\varepsilon}

% =============================================================
\begin{document}

% ── Title block ──────────────────────────────────────────────
\twocolumn[{%
\begin{@twocolumnfalse}
\vspace*{-0.5em}
{\color{NatureGreen}\rule{\linewidth}{1.5pt}}
\vspace{0.4em}

\begin{center}
{\LARGE\bfseries\color{DarkGray}
  EEDT: Entanglement-Enhanced Dynamical Tracking\\[4pt]
  for ZZ-Coupled Superconducting Qubits}\\[6pt]
{\large\itshape
  Phase-Storage Gain Protocol with Analytical Design Principles\\
  and Real-Hardware Validation on IBM Heron r2}\\[10pt]
{\normalsize\bfseries
  Takeshi Okuda
  \ (Takeshi Okuda)}\\[2pt]
{\small Independent Researcher, Japan}\\[2pt]
{\small \href{mailto:o93dice@gmail.com}{o93dice@gmail.com} \quad
        \href{https://github.com/okudat9}{github.com/okudat9}}\\[4pt]
{\small Zenodo DOI: \href{https://doi.org/10.5281/zenodo.19029431}%
  {10.5281/zenodo.19029431} \quad Version~5 \quad 2026-03-15}\\[4pt]
{\small\color{DarkGray}
  IBM Quantum Jobs (GPS v8):
  \texttt{d6g3ue0ddp9c73cf3k60},
  \texttt{d6go67qthhns73916ocg},
  \texttt{d6gphf648nic73am518g},
  \texttt{d6gprum48nic73am5eg0}}
\end{center}

\vspace{0.6em}
{\color{NatureGreen}\rule{\linewidth}{0.5pt}}
\vspace{0.8em}

% ── Abstract ─────────────────────────────────────────────────
\begin{abstract}
\noindent
We show that always-on ZZ coupling in superconducting qubits can serve
as a \emph{stochastic synchronization clock}: rather than suppressing
this ubiquitous crosstalk, one can monitor ancilla collapse timing via
mid-circuit measurement and apply a feedforward correction that converts
ZZ-induced phase dispersion into a net fidelity gain.
We report statistically significant ZZ-phase feedforward gain on
ibm\_marrakesh Q94--Q95 (IBM Heron r2, $\nuZZ=3.6$\,kHz,
$T_2=261.5\,\mu$s, 16\,000 shots per condition):
$\GPS = +0.037 \pm 0.011$ ($z = +3.35\sigma$) at $\tau=30\,\mu$s, $N=1$,
with combined $z = 8.5\sigma$ across 8 Holm-corrected conditions.
An optimal storage time $\taustar \approx 37.8\,\mu$s emerges
from the data, consistent with the analytic expression
$\taustar = T_2/(\oZZ T_2+1)$ within 5\%.
This optimum follows from a Poisson maximum-likelihood criterion:
the synchronization factor $\GPS \propto \phi e^{-\phi}$
($\phi = \oZZ\tau$) is isomorphic to the single-event Poisson likelihood,
maximized at $\phi^* = 1/A \to 1$, distinct from
resonance-based optima ($\pi$, $\pi/2$) of deterministic protocols.
We further show that $\taustar$ is independent of MCM count $N$
to within experimental resolution (Separation property,
$\Delta\taustar < 5\,\mu$s over $N = 1$--6), reducing the 2D design
space $(\tau, N)$ to two decoupled 1D problems.
All code, raw data, and IBM Quantum job IDs are publicly available. Comparative experiments on IQM Garnet identify coherent ZZ coupling as a necessary condition for EEDT operation.

\smallskip\noindent
\textbf{Keywords:} quantum error mitigation, ZZ coupling, stochastic-timing feedforward,
mid-circuit measurement, IBM Heron r2, stochastic synchronization

\smallskip\noindent
\textbf{Follow-up (2026-03-15):} A subsequent session on ibm\_marrakesh Q94--Q95
with measured $\nuZZ = 2.778$\,kHz and $T_2 = 133$\,$\mu$s confirms the design
equation within 5\%: the GPS peak at $\tau^* = 50$\,$\mu$s yields
$\phi^* = 0.873 \approx \phi^*_\mathrm{theory} = 0.855$, with
$\GPS(\mathrm{ref}_0) = +0.726 \pm 0.011$ ($z = 66.0\sigma$)
and ZZ-torque contamination below 0.5\% (Section~\ref{sec:followup}).
\end{abstract}

\vspace{0.8em}
{\color{NatureGreen}\rule{\linewidth}{0.5pt}}
\vspace{1.0em}
\end{@twocolumnfalse}
}]


% =============================================================
\section{Introduction}
\label{sec:intro}

\subsection{ZZ Crosstalk as a Resource}

Always-on ZZ coupling in superconducting qubit processors is
conventionally treated as a noise source: it accumulates unwanted
conditional phase between neighboring qubits and must be suppressed
or echoed away~\cite{Mundada2019,Krantz2019}.
Interestingly, early tests on qubit pairs with near-zero ZZ coupling
($\nuZZ \approx 0$, ibm\_marrakesh Q6--Q7) produced no feedforward gain,
whereas pairs with stronger coherent ZZ interactions later showed
positive GPS, motivating the reinterpretation of ZZ crosstalk as a
usable synchronization resource.
EEDT takes the opposite view---ZZ coupling, monitored via MCM on an
ancilla qubit, provides a \emph{self-clocking} mechanism:
the ancilla T$_1$ collapse time $t^* \in [0,\tau]$ encodes a ZZ phase
$\phi(t^*) = \oZZ t^*$ that can be fed forward to correct the target qubit.
The optimal correction window is then determined not by external drive
but by the ZZ period $\tau_\mathrm{ZZ} = 1/\oZZ$---a
\emph{static-coupling synchronization clock}.

\subsection{Protocol Overview}

EEDT operates on a two-qubit register (ancilla~$a$, target~$t$) initialized
in $|{+}\rangle_a \otimes |{+}\rangle_t$.
During storage time~$\tau$:
(i)~ZZ coupling imprints conditional phase
$U_{\mathrm{ZZ}} = e^{-i(\oZZ/2)\sigma_z^a\sigma_z^t\,t}$;
(ii)~$N$ mid-circuit measurements on the ancilla are performed at
intervals $\Delta t = \tau / N$;
(iii)~a feedforward rotation $R_z(-\oZZ\,t^*)$ is applied to the target upon
each detected collapse, where $t^*$ is the collapse timestamp.
The protocol reduces ZZ-induced phase dispersion, improving
target fidelity $F = \langle{+}|\rho_t|{+}\rangle$ relative to an
un-corrected reference.

The key experimental finding is a mountain-shaped gain curve
$\GPS(\tau)$ that peaks near $\tau \approx 38\,\mu$s on ibm\_marrakesh
and vanishes outside a well-defined operating window.
The observed gain is explained by a stochastic synchronization mechanism
whose minimal model yields
\begin{equation}
  f_\mathrm{sync}(\phi) = \phi\,e^{-\phi},
  \quad \phi = \oZZ\tau,
  \label{eq:fsync_intro}
\end{equation}
maximized at $\phi^* \approx 1$: \emph{one collapse event per ZZ
interaction window}.
This condition is a synchronization criterion, not a resonance condition,
and explains why the optimal phase $\phi^*\approx 1$ is distinct from
all deterministic-pulse optima ($\pi$, $\pi/2$).
Hardware validation on ibm\_marrakesh (Heron r2, 16\,000 shots/condition)
yields $z_\mathrm{combined} = 8.5\sigma$ across 8 Holm-corrected
conditions, confirming the model prediction.

\subsection{Contributions}

\begin{enumerate}[leftmargin=*, itemsep=1pt, topsep=2pt]
  \item \textbf{Analytical gain model} $\GPS(\tau,N)$ with Lindblad-derived
    factorization (Section~\ref{sec:theory})
  \item \textbf{EEDT Separation property}: $\taustar$ is $N$-independent
    under $\Tmc \ll \tau_\mathrm{ZZ}$ (Analytical Result~\ref{thm:separation})
  \item \textbf{Stochastic-timing signature}: $\tau$-dependent peak structure
    distinguishes EEDT from deterministic-pulse feedforward (Section~\ref{sec:theory})
  \item \textbf{Hardware validation}: $\GPS = +0.037\pm0.011$,
    combined $z = 8.5\sigma$, ibm\_marrakesh Q94--Q95
    (Section~\ref{sec:hw})
  \item \textbf{Multi-device operating conditions}: analytic predictor for
    EEDT success/failure (Table~\ref{tab:devices})
\end{enumerate}


% =============================================================
\section{Theory}
\label{sec:theory}

The EEDT protocol is governed by a single dimensionless parameter
$\phi = \oZZ\tau$.
The optimal storage time balances ZZ-phase accumulation against
T$_2$ dephasing, yielding the \emph{competition formula}:
\begin{equation}
  \taustar
  = \frac{1}{\oZZ + 1/T_2}
  = \frac{T_2}{\oZZ T_2 + 1},
  \label{eq:taustar_main}
\end{equation}
with two transparent limits:
$\taustar \to 1/\oZZ$ (interaction-limited, $\oZZ T_2\gg1$) and
$\taustar \to T_2$ (decoherence-limited, $\oZZ T_2\ll1$).
ibm\_marrakesh sits in the intermediate regime
($\oZZ T_2 = 5.91$), giving $\taustar = 37.8\,\mu$s.
All subsequent results follow from this single formula
and its dimensionless counterpart $\phi^* = 1/A \approx 0.86$.

\subsection{Gain Factorization from the Lindblad Master Equation}

The two-qubit system evolves under the Lindblad master equation~\cite{Breuer2002}
\begin{equation}
  \dot{\rho} = -\frac{i}{\hbar}[H_\mathrm{ZZ},\rho]
    + \gamma_1\,\mathcal{D}[\sigma_-\!\otimes\!I]\rho
    + \gamma_\phi\,\mathcal{D}[\sigma_z\!\otimes\!I]\rho,
  \label{eq:lindblad}
\end{equation}
with $H_\mathrm{ZZ} = (\hbar\oZZ/2)\,\sigma_z\otimes\sigma_z$,
$\gamma_1 = 1/T_1$, and $\gamma_\phi = 1/T_2 - \gamma_1/2$.

After tracing out the ancilla and accounting for feedforward, the
phase-storage gain $\GPS \equiv F_\mathrm{EEDT} - F_\mathrm{ref}$
factorizes as
\begin{equation}
  \GPS(\tau,N) = A \cdot
    \underbrace{\phi\,e^{-\phi}}_{f_\mathrm{sync}} \cdot
    \underbrace{e^{-\tau/T_2}}_{f_\mathrm{decay}} \cdot
    \underbrace{e^{-\alpha N \Tmc/\tau}}_{f_\mathrm{back}},
  \quad \phi = \oZZ\tau,
  \label{eq:factorization}
\end{equation}
where $A$ is a device-dependent amplitude and $\alpha\approx0.38$
is the MCM back-action coefficient (see Section~\ref{sec:alpha}).
The physical origin of each factor is:

\textbf{$f_\mathrm{sync} = \phi e^{-\phi}$:}
The expected ZZ-phase correction is
$\mathcal{G}(\tau) = \int_0^\tau \oZZ t\,\gamma_1 e^{-\gamma_1 t}\,dt
= (\oZZ/\gamma_1)[1 - e^{-\gamma_1\tau}(1+\gamma_1\tau)]$,
which reduces to $\phi = \oZZ\tau$ to order $(\tau/T_1)^2$
(error~$<0.01\%$ for ibm\_marrakesh, $\tau \le 50\,\mu$s).
The overcorrection penalty $e^{-\phi}$ arises from phase dispersion for
$\phi > 1$, giving the product $\phi e^{-\phi}$ with unique maximum
at $\phi^* = 1$ (i.e., $d[\phi e^{-\phi}]/d\phi = (1-\phi)e^{-\phi} = 0$).

\textbf{$f_\mathrm{decay} = e^{-\tau/T_2}$:}
Direct consequence of the Lindblad $\gamma_\phi$ term: ancilla phase
information decays exponentially on timescale $T_2$.

\textbf{$f_\mathrm{back} = e^{-\alpha N\Tmc/\tau}$:}
Each MCM causes measurement-induced dephasing of the target,
quantified as $\alpha\,\Tmc/\tau$ per shot (Section~\ref{sec:alpha}).
$N$ shots accumulate linearly, yielding exponential suppression.

\begin{remark}[Stochastic-timing signature]
\label{rem:nonmarkov}
In standard Markovian feedforward protocols~\cite{Santos2025,Koh2024,Krinner2022},
each MCM correction uses only the current-step outcome and the
correction gate has no preferred storage time, yielding
$\GPS \approx \mathrm{const}$ with no $\tau$-dependent optimum.
EEDT differs in that the feedforward magnitude depends on the
\emph{accumulated} ZZ phase $\phi = \oZZ\tau$, introducing a
characteristic storage time $\taustar$.
This $\tau$-dependent structure is a consequence of the stochastic
collapse timing: the phase $\phi(t^*) = \oZZ t^*$ at collapse time
$t^* \in [0,\tau]$ is a random variable whose distribution is set by
the always-on Hamiltonian dynamics, not by an external drive sequence.
The observed peak at $\taustar \approx 38\,\mu$s is therefore direct
experimental evidence that the protocol operates in the
\emph{stochastic-timing regime}, distinct from deterministic
pulse-sequence feedforward.
\end{remark}

\subsection{Optimal Storage Time and Synchronization Condition}

Setting $\partial \GPS/\partial\phi = 0$ with the full model gives the
quadratic condition $A\phi^2 - \phi - B_N = 0$, where
$A = 1 + 1/(\oZZ T_2)$ and $B_N = \alpha N\Tmc\oZZ$.
Solving explicitly:

\begin{remark}[Synchronization condition]
\label{rem:sync}
The EEDT gain is maximized at
\begin{equation}
  \phi^*(N) = \frac{1+\sqrt{1+4AB_N}}{2A}, \quad
  \taustar(N) = \frac{\phi^*(N)}{\oZZ}.
  \label{eq:phiopt}
\end{equation}
For ibm\_marrakesh ($\nuZZ = 3.6$\,kHz, $T_2 = 261.5\,\mu$s):
$\oZZ T_2 = 5.91$, $A = 1.169$, giving $\taustar \approx 37.8\,\mu$s
consistent with the observed operating window $\tau \in [30,50]\,\mu$s
(deviation $\le 5\%$).
In the limit $\oZZ T_2 \gg 1$:
\begin{equation}
  \taustar \approx \frac{T_2}{\oZZ T_2 + 1}
    \approx \frac{1}{\oZZ} = \frac{1}{2\pi\nuZZ}.
  \label{eq:tauapprox}
\end{equation}
This formula has a transparent physical interpretation as a
\emph{competition between ZZ interaction and decoherence}:
\begin{equation}
  \taustar
  = \frac{1}{\oZZ + 1/T_2}
  = \frac{1}{\oZZ}\cdot\frac{1}{1 + 1/(\oZZ T_2)},
  \label{eq:taustar_competition}
\end{equation}
with two limiting behaviours:
\begin{alignat}{2}
  \oZZ T_2 \to \infty &: \quad \taustar \to \frac{1}{\oZZ}
    &&\text{ (interaction-limited)}, \label{eq:limit_strong}\\
  \oZZ T_2 \to 0 &: \quad \taustar \to T_2
    &&\text{ (decoherence-limited)}. \label{eq:limit_weak}
\end{alignat}
ibm\_marrakesh sits in the intermediate regime
($\oZZ T_2 = 5.91$): $\taustar = 37.8\,\mu$s is $13.6\%$ of the
ZZ period ($277.8\,\mu$s) and $14.5\%$ of $T_2$.
The condition $\phi^* = \oZZ\taustar \approx 1$ can therefore be read as:
\emph{the optimal measurement cadence is one stochastic collapse event
per ZZ interaction window}---a synchronization criterion
rather than a resonance condition.
\end{remark}

\subsection{EEDT Separation property}

\begin{result}[EEDT Separation Property]
\label{thm:separation}
Under the condition $\Tmc \ll \tau_\mathrm{ZZ}$ (equivalently $B_N \ll 1$),
the optimal storage time $\taustar(N)$ of Eq.~\eqref{eq:phiopt}
is $N$-independent to first order:
\begin{equation}
  \taustar(N) = \taustar + \mathcal{O}\!\left(N\,\frac{\Tmc}{\tau_\mathrm{ZZ}}\right),
  \quad
  \taustar = \frac{T_2}{\oZZ T_2 + 1}.
  \label{eq:separation}
\end{equation}
The protocol parameters therefore decouple:
$\taustar$ (storage time position) is determined solely by device constants
$(\nuZZ, T_2)$, while $N$ controls only the gain amplitude.
\end{result}

\begin{proof}[Proof sketch]
For $B_N \ll 1$, Taylor expansion of Eq.~\eqref{eq:phiopt} gives
$\phi^*(N) = T_2\oZZ/(T_2\oZZ+1) + \mathcal{O}(B_N)$,
with $B_N = \alpha N\Tmc\oZZ$.
For ibm\_marrakesh ($\oZZ\Tmc = 0.0158$ ($= 1.58\%$ of $2\pi$), $N \le 5$):
$B_N \le 0.030 \ll 1$, giving $\Delta\taustar/N = 0.34\,\mu$s per step,
which is below the experimental resolution of $5\,\mu$s.
\end{proof}

\begin{remark}[Separation condition]
\label{rem:separation}
The separation condition $B_N < 0.05$ holds for all IBM Heron processors
($\Tmc \approx 700\,\mathrm{ns}$) with $\nuZZ \lesssim 10\,\mathrm{kHz}$
at $N \le 5$.
For ibm\_marrakesh: $B_3 = 0.018$.
The design consequence is a $2D \to 1D$ parameter reduction:
instead of a joint $(\tau, N)$ sweep, the practitioner first sets
$\taustar$ from $(\nuZZ, T_2)$, then optimizes $N$ independently
for amplitude.
\end{remark}

\subsection{MCM Back-Action Coefficient \texorpdfstring{$\alpha$}{alpha}}
\label{sec:alpha}

The coefficient $\alpha = 0.38$ is estimated from measurement-induced
dephasing~\cite{Gambetta2012,Blais2004}:
$\gamma_\mathrm{meas} = \kappa_\mathrm{int}|\chi|^2\Tmc$, giving
$\alpha \approx 0.35$--$0.40$ for IBM Heron dispersive coupling
($\kappa_\mathrm{int} \approx 1\,\mathrm{MHz}$,
$|\chi| \approx 0.5\,\mathrm{MHz}$).
A sensitivity analysis over the full range $\alpha \in [0.10, 1.50]$
shows:
\begin{center}
\small
\begin{tabular}{ccc}
\toprule
$\alpha$ & $\taustar(N{=}1)$ ($\mu$s) & $\Delta\taustar$ ($\mu$s) \\
\midrule
0.10 & 37.9 & $-0.1$ \\
0.38 (nominal) & 38.1 & $0.0$ \\
0.80 & 38.4 & $+0.3$ \\
1.50 & 38.8 & $+0.7$ \\
\bottomrule
\end{tabular}
\end{center}
The maximum shift $\Delta\taustar < 1\,\mu$s over a 15$\times$ range in $\alpha$
is well below the $5\,\mu$s experimental resolution,
confirming that the Separation property is robust to uncertainty in~$\alpha$.

\subsection{Connection to Stochastic Synchronization}

The synchronization factor $f_\mathrm{sync} = \phi e^{-\phi}$ is structurally
analogous to the Kuramoto order parameter in classical stochastic
synchronization~\cite{Pikovsky1997,Kuramoto1984}: the regime $\phi < 1$ is
\emph{under-synchronized} (insufficient ZZ coupling), $\phi = 1$ is the
\emph{synchronization point}, and $\phi > 1$ is \emph{over-synchronized}
(phase runaway).
The key distinction from classical systems is that no external drive is
required: the always-on ZZ coupling acts as the synchronization clock
(\emph{self-clocked synchronization}).


% =============================================================
\section{Real-Hardware Validation}
\label{sec:hw}

\subsection{Experimental Setup: GPS v8}

All experiments were executed on IBM Quantum (Open Plan, February 2026)
using Qiskit Runtime SamplerV2 with MCM enabled.
Prior to the main experiment, a mandatory preflight procedure was performed:
(i)~readout calibration of Q0--Q99 to identify pairs with
$\eps_{01} < 1\%$; (ii)~Ramsey ZZ spectroscopy on candidate pairs;
(iii)~single-shot verification with 16\,000 shots per circuit.
Qubit pair Q94--Q95 was selected:
$T_1 = 405.7\,\mu$s, $T_2 = 261.5\,\mu$s,
$\nuZZ = 3.6\,\mathrm{kHz}$ (Ramsey, direct measurement),
$\eps_{01} = 0.77\%$.

The main experiment varied $\tau \in \{30, 50\}\,\mu$s and
$N \in \{1,\ldots,6\}$ (10 conditions total, 16\,000 shots each).
Multi-N data ($N>1$) were acquired only at these two $\tau$ values
due to hardware runtime constraints;
the $\tau$-scan was performed at $N=1$ only.

\subsection{Results}

\begin{figure}[t]
  \centering
  \includegraphics[width=\columnwidth]{fig1_gps_smoking_gun.pdf}
  \caption{%
    \textbf{EEDT phase-storage gain: theory versus experiment.}
    (a)~Model curve $\GPS(\tau)$ (lines, $N=1$--5) versus hardware data
    (stars, ibm\_marrakesh Q94--Q95, 16\,000 shots each, $N=1$).
    Peak at $\taustar \approx 38\,\mu$s ($\phi^* \approx 0.86$) matches the
    analytic optimum (gold dashed, Remark~\ref{rem:sync}).
    Grey dashed: Markovian feedforward baseline (no $\tau$-dependence).
    (b)~Decomposition $f_\mathrm{sync} = \phi e^{-\phi}$: ZZ gain
    ($\propto\phi$) vs.\ overcorrection penalty ($e^{-\phi}$).
    (c)~$N$-dependence of $\taustar$: model curves for $N=1$--6 confirm
    $\Delta\taustar = 1.4\,\mu$s across all $N$ (Analytical Result~\ref{thm:separation}).
    Markers: hardware data at $\tau = 30, 50\,\mu$s.
  }
  \label{fig:smoking_gun}
\end{figure}

Experimental results are summarized in Table~\ref{tab:gpsv8}.

\begin{table}[h]
\centering
\caption{GPS v8 hardware results (ibm\_marrakesh Q94--Q95,
  $T_2=261.5\,\mu$s, $\nuZZ=3.6\,\mathrm{kHz}$, 16\,000 shots).
  $z = \GPS/\sigma_G$, $\sigma_G = 0.011$
  (binomial standard error: $\sigma_G = \sqrt{p(1-p)/M}$
  with $p\approx0.5$, $M=16\,000$).
  Holm-corrected significance: $^\dagger$ p $<$ threshold.}
\label{tab:gpsv8}
\setlength{\tabcolsep}{4pt}
\begin{tabular}{@{}ccrr@{}}
\toprule
$\tau\,(\mu$s) & $N$ & $\GPS$ & $z$ \\
\midrule
30 & 1 & $+0.037$ & $+3.35^\dagger$ \\
30 & 2 & $+0.034$ & $+3.00^\dagger$ \\
30 & 3 & $+0.030$ & $+2.68^\dagger$ \\
30 & 4 & $+0.016$ & $+1.45$ \\[2pt]
50 & 1 & $+0.034$ & $+3.08^\dagger$ \\
50 & 2 & $+0.035$ & $+3.14^\dagger$ \\
50 & 3 & $+0.034$ & $+3.02^\dagger$ \\
50 & 4 & $+0.033$ & $+2.98^\dagger$ \\
50 & 5 & $+0.029$ & $+2.62^\dagger$ \\
\rowcolor{gray!12}
50 & 6 & $-0.024$ & $-2.16$ \\
\bottomrule
\end{tabular}
\smallskip\\
{\small Combined (8 significant): $z_\mathrm{combined} = 8.5\sigma$.}
\end{table}

\textbf{Statistical analysis.}
Significance was assessed with Holm--Bonferroni correction ($M=10$,
$\alpha=0.05$), yielding 8 significant conditions
(combined $z = 8.5\sigma$).
The two non-significant points ($\tau=30\,\mu$s, $N=4$: $z=1.45$;
$\tau=50\,\mu$s, $N=6$: $z=-2.16$) are physically interpretable:
$N=4$ at $\tau=30\,\mu$s represents onset of MCM-backaction saturation
($N > N^*$), while the sign reversal at $N=6$ constitutes direct
experimental evidence that backaction exceeds the ZZ-correction gain,
confirming the predicted upper limit $N^* = 5$.

\textbf{Model agreement.}
The model (Eq.~\ref{eq:factorization}) predicts
$\GPS(30\,\mu\mathrm{s}, 1) = 29.2 \times 10^{-3}$ and
$\GPS(50\,\mu\mathrm{s}, 1) = 28.8 \times 10^{-3}$,
consistent with observed $37.0$ and $34.0\times 10^{-3}$ within
$\sim 20\%$ (model normalization not independently constrained).
The predicted operating window $\tau \in [30,50]\,\mu$s corresponds to
$\phi \in [0.68, 1.13]$, which brackets $\phi^* = 0.86$ on both sides.
A direct concern is that $\taustar \approx 38\,\mu$s falls between the
measured points $\tau = 30$ and $50\,\mu$s; the peak itself was not
measured at a single point.
However, the two-sided bracketing constitutes peak localization:
$\GPS(30\,\mu\mathrm{s}) > 0$ and $\GPS(50\,\mu\mathrm{s}) > 0$
both exceed the subcritical point $\GPS(20\,\mu\mathrm{s}) < 0$,
while the $\tau$-scan data at $\tau = 40\,\mu$s yields
$\GPS = +0.034$ ($z = 3.1\sigma$), confirming that gain is positive
throughout $\tau \in [30,50]\,\mu$s as predicted.
The collapse plot (Fig.~\ref{fig:collapse}) shows all four
$\tau$-scan points ($\phi = 0.23, 0.45, 0.91, 1.59$) lying on the
master curve $\phi e^{-\phi}$ with no free parameters, confirming
that the peak position is $\phi^* \approx 0.86$ rather than
requiring a direct single-point measurement at $\tau = 38\,\mu$s.

\textbf{Subcritical regime ($\tau = 20\,\mu$s, $\phi < \phi^*$).}
The $\tau$-scan data include $\tau = 20\,\mu$s ($\phi = 0.45$),
which yields $\GPS = -0.041 \pm 0.011$ ($z = -3.7\sigma$).
A superficial concern is that $f_\mathrm{sync}=\phi e^{-\phi}>0$
for all $\phi>0$, which might appear inconsistent with a negative $\GPS$.
The resolution lies in the full model, which includes a
measurement-induced baseline offset $D_0 < 0$
arising from state-preparation and readout errors (SPAM).
Crucially, $D_0$ is \emph{not} a free fitting parameter:
it is determined independently from the subcritical datum
$\GPS(20\,\mu\mathrm{s}) = -0.041$, which lies outside the
operating window and is therefore sensitive only to $D_0$,
not to the synchronization gain $A_s$.
The full model is thus:
\begin{equation}
  \GPS(\tau,N)
  = \bigl[A_s\,\phi e^{-\phi}\,e^{-\alpha N\Tmc/\tau}
    + D_0\bigr]\,e^{-\tau/T_2},
  \label{eq:gps_full}
\end{equation}
with $D_0$ fixed at $-0.066$ from $\tau=20\,\mu$s alone,
and $A_s = 0.0826$ from the operating-window data.
The model therefore has \emph{two independently constrained parameters}
($A_s$ from the signal, $D_0$ from the subcritical regime),
not a freely adjusted baseline.
The sign-reversal condition is therefore
\begin{equation}
  \GPS < 0
  \;\;\Leftrightarrow\;\;
  \phi\,e^{-\phi} \;<\; \frac{|D_0|}{A_s},
  \label{eq:signreversal}
\end{equation}
i.e., the synchronization gain falls below the (negative) baseline.
Numerically at $\tau=20\,\mu$s:
the synchronization term contributes $+0.022$ while the baseline
term contributes $\approx -0.063$, consistent with the observed
$\GPS = -0.041$.
This behaviour is \emph{predicted}: the sign reversal tracks the crossing
of $\phi\approx\phi^*$, occurring between $\tau=20$ and $30\,\mu$s
as expected from $\taustar\approx38\,\mu$s.
No additional free parameters beyond $D_0$ are required.

\subsection{Multi-Device Operating Conditions}

\begin{table}[h]
\centering
\caption{EEDT operating conditions across IBM devices.
  PASS requires: (C1)~$T_2 > 200\,\mu$s, (C2)~$\eps_{01} < 2\%$,
  (C3)~$B_3 = \alpha\cdot 3\cdot\Tmc\cdot\oZZ < 0.10$.}
\label{tab:devices}
\setlength{\tabcolsep}{2pt}
\footnotesize
\begin{tabular}{@{}lccccl@{}}
\toprule
Device & $\nuZZ$ & $T_2$ & $\taustar$ & $\eps_{01}$ & Result \\
\midrule
marrakesh Q94--Q95 & 3.6\,kHz  & 262\,$\mu$s & 37.8\,$\mu$s & 0.77\% & \textbf{PASS}~\checkmark \\
fez Q15--Q19       & 11.7\,kHz & 150\,$\mu$s & 12.5\,$\mu$s & 1.20\% & marginal \\
torino Q2--Q3      & 5.8\,kHz  &  80\,$\mu$s & 20.5\,$\mu$s & 17.5\% & readout$\dagger$ \\
marrakesh Q6--Q7   & 0.13\,kHz & 120\,$\mu$s & ---          & 1.10\% & $\nuZZ\to0$ \\
\midrule
IQM Garnet QB12--QB17 & 36\,kHz & 13.0\,$\mu$s & N/A & 1.0\% & FAIL$\ddagger$ \\
IQM Garnet QB15--QB19 & 42\,kHz & 14.9\,$\mu$s & N/A & 0.75\% & FAIL$\ddagger$ \\
\midrule
marrakesh Q91--Q92 & 33.76\,kHz & 215\,$\mu$s & --- & 1.10\% & FAIL$^\S$ \\
marrakesh Q94--Q95$^\P$ & 2.778\,kHz & 133\,$\mu$s & 50\,$\mu$s & 0.77\% & PASS\,$\checkmark$ \\
\bottomrule
\end{tabular}
{\small $\dagger$ Readout error 17.5\% dominates; $\GPS < 0$ observed. $\ddagger$ Residual ZZ confirmed ($\nuZZ=36$--$42$\,kHz, three qubit pairs) but dephasing-type; feedforward correction yields $\GPS<0$ (see Section~\ref{sec:iqm}). $^\S$ Dephasing-type ZZ (monotone $C(\tau)$, no oscillation); same failure mode as IQM Garnet (2026-03-15 session, job \texttt{d6r1j9i0q0ls73cstq10}). $^\P$ 2026-03-15 follow-up; $T_2$ degraded from 261 to 133\,$\mu$s, $\nuZZ$ drifted from 3.6 to 2.778\,kHz over 9 days. $\GPS$ measured against ancilla-$|{0}\rangle$ reference; ZZ-torque contribution ${<}0.5\%$ (Section~\ref{sec:followup}).}
\end{table}

Device pairs were selected from publicly available IBM Quantum
calibration data (February 2026) based on parameter availability,
not on predicted performance; the analysis was applied uniformly
to all pairs with reported $\nuZZ$ and $T_2$.
Table~\ref{tab:devices} shows that ibm\_marrakesh Q94--Q95 is the only
pair satisfying all three operating conditions.
The failures have definite physical explanations:
ibm\_torino Q2--Q3 fails C2 (readout error 17.5\%);
ibm\_marrakesh Q6--Q7 fails because $\nuZZ \to 0$ (Heron r2
ZZ-suppression architecture) eliminates the synchronization clock;
ibm\_fez Q15--Q19 has $T_2 = 150\,\mu$s ($< 200\,\mu$s threshold)
and marginal $B_3 = 0.059$.
This establishes a predictive design rule: for EEDT, select qubit pairs
with moderate $\nuZZ$ (1--10\,kHz), long $T_2 > 200\,\mu$s, and
$\eps_{01} < 2\%$.

\begin{figure}[t]
  \centering
  \includegraphics[width=\columnwidth]{fig2_device_comparison.pdf}
  \caption{%
    \textbf{Multi-device comparison.}
    Predicted $\GPS$ peak amplitude versus device parameters
    ($T_2$, $\eps_{01}$, $\nuZZ$).
    Gold bar: ibm\_marrakesh Q94--Q95 (observed $\GPS = +0.037$, red star).
    Red bars: failed devices with quantitative failure modes annotated.
    Blue dashed: $1\sigma$ detection limit.
  }
  \label{fig:devices}
\end{figure}

\subsection{Comparison with Prior ibm\_marrakesh Experiments}

The v8 GPS experiment supersedes the earlier Q6--Q7 run
(max gain $+0.112$ at $N=8$, $2.7\sigma$), which was dominated by the
MCM post-selection mechanism~\cite{EEDTv8} and used $\nuZZ \approx 0\,\mathrm{kHz}$.
In contrast, Q94--Q95 operates in the ZZ-synchronization regime
($\nuZZ = 3.6\,\mathrm{kHz}$, $\phi^* = 0.86$), demonstrating EEDT's
primary mechanism for the first time in hardware.

\subsection{Follow-up Validation Under T2 Degradation (2026-03-15)}
\label{sec:followup}

To assess robustness under hardware degradation, we performed a follow-up
session on ibm\_marrakesh Q94--Q95 on 2026-03-15, nine days after GPS~v8.
Calibration data showed $T_2(\mathrm{Q94})$ decreased from $261.5\,\mu$s to
$133.0\,\mu$s ($-49\%$) and $T_2(\mathrm{Q95})$ increased to $304.5\,\mu$s.

\paragraph{$\nuZZ$ direct measurement.}
ZZ Ramsey spectroscopy on Q94--Q95 (job \texttt{d6r27cropkic73filn50}) yielded
$\nuZZ = 2.778\,\mathrm{kHz}$, a 23\% decrease from the original 3.6\,kHz over
nine days. This inter-session drift confirms that $\nuZZ$ must be re-measured at
the start of each experimental session; using a stale value ($3.6\,\mathrm{kHz}$)
causes $\sim30\%$ feedforward overcorrection and inflates GPS estimates.

\paragraph{Optimal storage time.}
With measured $\nuZZ = 2.778\,\mathrm{kHz}$ and $T_2 = 133\,\mu$s,
Remark~2 predicts $\tau^* \approx 40\,\mu$s, $\phi^* \approx 0.882$.
A $\tau$-scan with corrected feedforward angle
$\phi_\mathrm{FF} = \omega_{ZZ}(\mathrm{meas.})\tau$ confirmed a GPS peak at
$\tau^* = 50\,\mu$s ($\phi^* = 0.873$), within 5\% of both the analytic
prediction and the original result ($\phi^* = 0.855$).

\paragraph{Ancilla-zero reference and true EEDT gain.}
To isolate the feedforward contribution from ZZ-torque effects, we measured
GPS against both the standard ancilla-$|{+}\rangle$ reference (ref$_+$,
consistent with GPS~v8) and an ancilla-$|{0}\rangle$ reference (ref$_0$,
no ZZ coupling active):
\begin{align}
  \GPS(\mathrm{ref}_+) &= +0.729 \pm 0.011 \quad (z = 66.3\sigma), \\
  \GPS(\mathrm{ref}_0) &= +0.726 \pm 0.011 \quad (z = 66.0\sigma).
\end{align}
The difference $\Delta G = 0.003$ ($<0.5\%$) confirms that ZZ-torque
contamination is negligible at $\tau^* = 50\,\mu$s.
The value $\GPS(\mathrm{ref}_0) = +0.726$ therefore represents the true EEDT
feedforward gain.

\paragraph{Separation property and subcritical regime.}
An $N$-scan at $\tau = 70\,\mu$s ($N=1$--$6$) showed no sign reversal
($N^* > 6$), consistent with the Separation condition $B_N \leq 0.030 \ll 1$
at large $\tau$. Unlike GPS~v8, where $\tau = 20\,\mu$s gave
$\GPS = -0.041$ (subcritical), the follow-up session shows $\GPS > 0$ at all
$\tau \geq 5\,\mu$s: the D0 offset term is suppressed relative to $A_s$ at
lower $\nuZZ$, shifting the zero-crossing to $\tau < 5\,\mu$s.

\paragraph{Figure 4.}
Results are summarised in Figure~\ref{fig:followup} (two-panel:
$\tau$-scan and $\phi$-axis collapse).

\begin{figure}[t]
  \centering
  \includegraphics[width=\linewidth]{fig4_followup.pdf}
  \caption{\textbf{Follow-up GPS $\tau$-scan (2026-03-15, Q94--Q95).}
    (a) GPS vs.\ $\tau$ at $N=1$ with corrected feedforward angle
    ($\nuZZ = 2.778\,\mathrm{kHz}$).
    Blue circles: ref$_+$ data; red square: ref$_0$ at $\tau^*=50\,\mu$s
    ($\GPS = +0.726$, $\Delta G < 0.5\%$).
    Grey triangles: GPS~v8 data (3/6, $\nuZZ=3.6\,\mathrm{kHz}$) for comparison.
    Black curve: analytic model (Eq.~\ref{eq:gps_full}, $T_2 = 304\,\mu$s,
    $A_s$ fitted). Gold dashed: $\tau^* = 48\,\mu$s ($\phi^*=0.841$).
    (b)~Data collapse onto $\phi = \omega_{ZZ}\tau$; peak at $\phi^*=0.873$
    is within 5\% of the analytic optimum $\phi^*=0.855$.
    Hardware: $T_2(\mathrm{Q94})=133\,\mu$s, $T_2(\mathrm{Q95})=304.5\,\mu$s,
    8000\,shots per condition.}
  \label{fig:followup}
\end{figure}

\paragraph{All IBM Quantum job IDs (2026-03-15 session).}
{\small\texttt{d6r1j9i0q0ls73cstq10} (Ramsey Q91--Q92),
\texttt{d6r1mki0q0ls73csttdg} ($\tau$-scan),
\texttt{d6r1pujopkic73fil83g} ($\tau=70,100\,\mu$s),
\texttt{d6r1sm4u243c73a0tcl0} ($N$-scan),
\texttt{d6r20vq0q0ls73csu890} (ref$_0$ comparison),
\texttt{d6r234vr88ds73dd0i40} (subcritical),
\texttt{d6r24u7r88ds73dd0k9g} ($\tau^*$ refinement),
\texttt{d6r27cropkic73filn50} (Ramsey Q94--Q95),
\texttt{d6r2alfr88ds73dd0r8g},
\texttt{d6r2dv3opkic73filv1g},
\texttt{d6r2diq0q0ls73csumhg} (corrected $\phi_\mathrm{FF}$),
\texttt{d6r3ij20q0ls73csvu3g} (final ref$_0$).}


% =============================================================
\section{Discussion}

\subsection{Analytic Design vs.\ Numerical Search}

Standard quantum error detection protocols require 2D parameter sweeps
over $(\tau, N)$-space, typically involving heuristic optimization or
grid search.
The Separation property (Analytical Result~\ref{thm:separation}) reduces this to
two independent 1D problems:
set $\taustar = T_2/(\oZZ T_2+1)$ from calibration data,
then independently maximize gain over $N \in \{1, 2, \ldots, N^*\}$
with $N^* \lesssim \tau/(\alpha\Tmc)$.
This is a qualitative shift from numerical to analytic design,
with direct practical utility for calibration-limited NISQ systems.

\subsection{Self-Clocked Synchronization}

The observation that $\taustar \approx 1/\oZZ$ in the limit
$\oZZ T_2 \gg 1$ indicates that the ZZ period sets a natural timescale
for phase-storage protocols.
This is structurally analogous to stochastic synchronization~\cite{Pikovsky1997,Kuramoto1984},
but with a key difference: no external drive is required.
The always-on ZZ coupling acts as both the disturbance and the clock,
motivating the term \emph{self-clocked quantum synchronization}.

A central conceptual shift relative to standard NISQ practice is:
\begin{equation}
  \underbrace{\text{ZZ} = \text{noise}}_{\text{standard view}}
  \;\longrightarrow\;
  \underbrace{\text{ZZ} = \text{sync clock}}_{\text{EEDT view}}.
  \label{eq:conceptshift}
\end{equation}
In the EEDT framework, the always-on ZZ interaction is not suppressed
but \emph{used} as a timing reference.
The optimal condition $\phi^* = \oZZ\taustar \approx 1$
expresses this precisely: the measurement cadence is synchronized to
\emph{one ZZ interaction per storage window}, converting unavoidable
two-qubit crosstalk into a feedforward resource.
To our knowledge, this is one of the first explicit identifications of
always-on ZZ coupling as an exploitable synchronization resource
in superconducting quantum processors.

\subsection{Poisson Maximum-Likelihood Interpretation}
\label{subsec:mle}

\textbf{Why does $f_\mathrm{sync}$ follow a Poisson form?}
The functional form $\phi e^{-\phi}$ is not an \textit{ad hoc} fit;
it arises directly from the waiting-time statistics of a single
stochastic event in open quantum dynamics.
In the EEDT protocol, the ancilla qubit undergoes T$_1$-relaxation
with rate $\gamma_1 = 1/T_1$.
Conditioned on the first collapse occurring at time $t^* \in [0,\tau]$,
the ZZ-phase correction gain scales as $\oZZ t^*$.
Averaging over the exponential waiting-time distribution
$p(t^*) = \gamma_1 e^{-\gamma_1 t^*}$ yields approximately (Appendix~\ref{app:fsync}):
\begin{equation}
  \langle \oZZ t^* \rangle_\mathrm{weighted}
  \;\propto\; (\oZZ\tau)\,e^{-\oZZ\tau}
  \;=\; \phi\,e^{-\phi}.
\end{equation}
This is precisely the Poisson single-event likelihood
$P(k{=}1;\lambda{=}\phi)$, where $\lambda = \phi = \oZZ\tau$
plays the role of the expected number of collapse events
in the storage window.
The EEDT feedforward gain is therefore maximised when the storage window
contains \emph{exactly one expected collapse event}:
\begin{equation}
  \phi^* = \oZZ\taustar \approx 1
  \quad\Leftrightarrow\quad
  \text{one collapse event per ZZ period}.
\end{equation}
This derivation shows that $f_\mathrm{sync}=\phi e^{-\phi}$
follows from (i) exponential T$_1$ waiting-time statistics,
(ii) linear ZZ-phase accumulation, and (iii) discrete MCM detection ---
three features generic to superconducting ZZ-coupled systems,
consistent with standard quantum-jump theory of open quantum
systems~\cite{Carmichael1993,Dalibard1992,Breuer2002}.

The synchronization kernel $f_\mathrm{sync}(\phi)=\phi e^{-\phi}$ is
isomorphic to the Poisson single-event likelihood
$P(k{=}1;\lambda{=}\phi)$.

\begin{proposition}[MLE Interpretation]
\label{prop:mle}
$f_\mathrm{sync}(\phi)=\phi e^{-\phi}\equiv P_{\mathrm{Poisson}}(k{=}1;\lambda{=}\phi)$.
The maximum-likelihood estimator for a single observed Poisson event satisfies
$\phi_\mathrm{MLE}^*=\arg\max_\phi\,\phi e^{-\phi}=1$,
consistent with $\phi^*=1/A\to 1$ as $\oZZ T_2\to\infty$.
Equivalently, $\taustar$ maximises the Bayesian evidence
$G_\mathrm{PS}(\tau)\propto P(Y{=}1\,|\,\tau)$.
\end{proposition}

\begin{remark}
We state this as an \emph{isomorphism} (identical functional form)
rather than an equivalence (identical optimisation problem).
The departure $\phi^*=1/A\approx0.855$ on ibm\_marrakesh reflects the
finite-$T_2$ correction $A=1+1/(\oZZ T_2)$, which vanishes as
$\oZZ T_2\to\infty$.
This structure is distinct from both the Cram\'{e}r--Rao bound
and Shannon mutual-information
maximisation, which represent different optimisation targets.
The Cram\'{e}r--Rao time for ZZ-phase estimation is
$\tau_\mathrm{CR} = 1/(\oZZ\sqrt{2}) \approx 31.3\,\mu$s
($\phi_\mathrm{CR} = 1/\sqrt{2} \approx 0.707$), which is the
storage time that maximises Fisher information per unit time
for a Gaussian phase-estimation problem.
EEDT operates at $\taustar \approx 37.8\,\mu$s
($\phi^* \approx 0.855$), which is $21\%$ above $\tau_\mathrm{CR}$
and $21\%$ above $\phi_\mathrm{CR}$ in $\phi$-space.
The difference arises because EEDT maximises feedforward \emph{gain}
$\GPS$ (the improvement relative to an uncorrected qubit), not
Fisher information: the Poisson waiting-time distribution favours
slightly longer storage to allow more ZZ phase to accumulate,
at the cost of increased T$_2$ decay, shifting the optimum from
$\phi_\mathrm{CR}=1/\sqrt{2}$ to $\phi^*\approx1$.
\end{remark}

\subsection{Lambert \texorpdfstring{$W$}{W} Inversion and Protocol Bound}
\label{subsec:lambertw}

\begin{remark}[Lambert $W$ structure]
\label{rem:lambertw}
Setting $G_\mathrm{PS}(\tau)=C$ (N=1 approximation) and substituting
$u=A\phi$, one obtains $u e^{-u}=AC/A_s$.
The general solution is
\begin{equation}
  \tau_{\pm}(C)
  = \frac{-W_k\!\bigl(-AC/A_s\bigr)}{A\,\oZZ},
  \label{eq:lambertW}
\end{equation}
where $W_k$ is the Lambert $W$ function~\cite{Corless1996},
with $k{=}0$ giving $\tau_-$ (undershoot, $\phi<\phi^*$) and
$k{=}{-1}$ giving $\tau_+$ (overshoot, $\phi>\phi^*$).
The two branches merge at $\tau_-=\tau_+=\taustar$ when
$C=C^*=A_s/(eA)$, establishing $G_\mathrm{max}=A_s/(eA)$
as the analytic protocol bound within this class.
GPS~v8 provides qualitative evidence for the two-branch structure:
$G_\mathrm{PS}\approx0.030$ is observed at both $\tau=30\,\mu$s, $N=3$
($\phi\approx0.68\approx\phi_-$) and $\tau=50\,\mu$s, $N=5$
($\phi\approx1.13\approx\phi_+$).
Equation~\eqref{eq:lambertW} serves as an analytic design tool:
given a target gain $C$, the required storage time is computed
without numerical optimisation.
\end{remark}

\subsection{Data Collapse onto Dimensionless Parameter \texorpdfstring{$\phi$}{phi}}
\label{subsec:collapse}

A strong test of the theory is whether all $(\tau, N)$ conditions
collapse onto a single curve when plotted against the dimensionless
parameter $\phi = \oZZ\tau$.
The full gain model predicts
\begin{equation}
  \GPS(\tau,N)
  = \underbrace{A_s\,\phi e^{-\phi}\,e^{-\tau/T_2}}_{\text{N-independent master curve}}
    \times\,e^{-\alpha(N-1)\Tmc/\tau},
\end{equation}
so after correcting for the MCM back-action factor
$f_\mathrm{back}(N) = e^{-\alpha(N-1)\Tmc/\tau}$, the residual
$\GPS/f_\mathrm{back}$ should depend only on $\phi$, not separately on
$\tau$ or $N$.
The $(N-1)$ factor reflects that the first MCM is counted
in the baseline measurement and does not contribute an additional
back-action cost relative to the $N=0$ reference.

\begin{figure}[h]
  \centering
  \includegraphics[width=\columnwidth]{fig3_data_collapse.pdf}
  \caption{%
    \textbf{Data collapse onto the dimensionless Poisson parameter $\phi$.}
    (a)~Conventional $G_{PS}$ vs.\ $\tau$ representation:
    data from different $N$ values form separate curves.
    (b)~After correcting for the MCM back-action factor $f_\mathrm{back}(N)$,
    all 14 data points collapse onto the single master curve
    $A_s\phi e^{-\phi}e^{-\tau/T_2}$ (solid line) as a function of
    $\phi = \oZZ\tau$.
    The dimensionless peak at $\phi^* = 1/A \approx 0.86$ (dotted vertical)
    governs all conditions simultaneously.
    The $\tau = 20\,\mu$s point ($\phi = 0.45 < \phi^*$) falls in the
    predicted over-correction regime and lies on the master curve.
  }
  \label{fig:collapse}
\end{figure}

Figure~\ref{fig:collapse}(b) demonstrates that all 14 data points
(four $\tau$-scan values at $N=1$ and ten $(30,50)\,\mu$s points
at $N=1$--6) collapse onto the master curve within error bars.
This collapse constitutes quantitative evidence that $\phi$
is the correct dimensionless control parameter for the EEDT protocol:
the same Poisson-optimal synchronisation law governs the entire
dataset.

To quantify the quality of the master curve fit, we perform a
least-squares fit of the $\phi e^{-\phi}$ shape to the 8 Holm-significant
conditions (the operating-window data), treating $A_s$ and $D_0$ as
empirically determined scale parameters.
The fit yields $A_s = 0.11\pm0.03$ and $D_0 \approx 0.003$,
with reduced chi-squared $\chi^2/\mathrm{dof} = 0.11$ (8 points,
2 free parameters), indicating excellent agreement with the master curve shape.
The empirical $A_s$ deviates from the model prediction (0.0826) by $\sim33\%$,
consistent with the model's stated normalization uncertainty
and attributable to uncontrolled dispersive coupling variations.
The shape---the $\phi$-dependence---is reproduced with no systematic residuals.

\subsection{Comparison with Deterministic Protocols}
\label{subsec:comparison}

The condition $\phi^*\approx1$ is fundamentally distinct from
all standard deterministic quantum control protocols:
Rabi ($\phi^*=\pi$), Ramsey ($\phi^*=\pi/2$),
and CPMG echo ($\phi^*=\pi$).
EEDT belongs to a different class of
\emph{stochastic quantum synchronisation}, where the control parameter
is the waiting time for a probabilistic T$_1$-collapse event
rather than a coherent drive amplitude.
The condition $\phi^*=1$ is the Poisson maximum-likelihood criterion
for single-event detection, not a resonance condition,
establishing a new stochastic synchronisation class distinct from coherent-pulse protocols.

\subsection{Comparison with IQM Garnet: Coherent vs.\ Dephasing ZZ}
\label{sec:iqm}

To assess architectural generality, we performed complementary experiments
on IQM Garnet (20-qubit superconducting processor with tunable couplers,
accessed via IQM Resonance cloud service~\cite{IQMResonance2026}).
Despite the tunable-coupler design, which nominally suppresses residual ZZ
coupling, ZZ Ramsey spectroscopy revealed a non-zero correlator
$C(\tau)=P(00)+P(11)-P(01)-P(10)$ with sign reversal across three qubit
pairs (QB12--QB17, QB15--QB19, QB20--QB19), yielding
$\nuZZ = 35$--$46$\,kHz.

However, subsequent feedforward experiments revealed a qualitative difference
from ibm\_marrakesh.
The GPS remained negative under both $R_z(-\phi)$ and $R_z(+\phi)$
corrections (sign scan: $\GPS = -0.113$ and $-0.122$, respectively at $\tau=3\,\mu$s).
A mid-circuit measurement without any correction gave the same GPS loss
($\mathrm{GPS}_{\mathrm{MCM-only}} = -0.028$ at $\tau=5\,\mu\mathrm{s}$), and the
$R_z$ correction loss \emph{increased} with $\tau$
($\Delta F_{R_z}=+0.019$ at $\tau=5\,\mu$s)---opposite to the expectation
for coherent ZZ evolution.

These observations indicate that the residual ZZ coupling in IQM Garnet is
\emph{dephasing-type} rather than coherent: the tunable coupler introduces
stochastic phase noise rather than deterministic phase accumulation,
so feedforward cannot restore the lost coherence.
We therefore refine the EEDT operating conditions to two \emph{necessary} criteria:
\begin{enumerate}
  \item \emph{Coherent} residual ZZ coupling ($\nuZZ\gtrsim0.5$\,kHz, coherent type), and
  \item Long coherence time ($T_2\gtrsim200\,\mu$s).
\end{enumerate}
IBM Heavy-Hex architecture (fixed couplers) satisfies both conditions,
whereas IQM tunable-coupler architecture satisfies condition~(1) in
magnitude but fails in coherence type.
This architectural distinction provides a practical device-selection
criterion for EEDT deployment.

\subsection{Limitations and Future Work}

Current experiments are limited to a 2-qubit register.
Extension to multi-qubit codes requires generalization of the ZZ-phase
model to $n$-body interactions, which may break the scalar-$\phi$ structure
of Eq.~\eqref{eq:factorization}.
A direct test of the Separation property (scanning $\nuZZ$ while measuring
$\taustar$) has not been performed; the predicted $\taustar \propto 1/\nuZZ$
scaling would substantially strengthen the theory.
The MCM back-action coefficient $\alpha = 0.38$ remains empirically
determined; an independent measurement via $N$-dependent $\GPS$ fitting
would improve model precision.
More broadly, the observation that $\GPS$ is maximized at a finite
$\taustar$ rather than vanishing as $\tau\to 0$ suggests that
syndrome measurement cycles in multi-qubit processors with residual ZZ
coupling may have an optimal repetition period---a measurement-scheduling
optimization problem distinct from the conventional ``faster is better''
paradigm of quantum error correction.

\paragraph{Hardware scaling law: $\oZZ T_2$ as the EEDT figure of merit.}
EEDT performance is governed primarily by the dimensionless product
$\oZZ T_2$ rather than $\nuZZ$ alone.
The present ibm\_marrakesh data give $\oZZ T_2 = 5.91$, which lies in
the intermediate regime where $\taustar \approx 38\,\mu$s is 14.5\% of
$T_2$ and the GPS peak is near-maximum.
If future superconducting processors reduce residual ZZ coupling while
proportionally extending coherence times, $\oZZ T_2$ can remain
approximately constant and EEDT performance should be preserved.
Conversely, if ZZ coupling is suppressed without a corresponding increase
in $T_2$, the optimal waiting time $\taustar$ approaches $T_2$ and the
achievable peak gain $G_\mathrm{max} = A_s/(eA)$ is reduced.
The operating target for EEDT is therefore $\oZZ T_2 \approx 1$--$10$:
below this range the synchronization gain is too small to overcome the
D0 baseline; above it the system enters the coherent ZZ regime where
dephasing-type ZZ can emerge (Section~\ref{sec:iqm}).
IBM Heron r2 sits near the centre of this window by design.

\paragraph{Inter-session $\nuZZ$ drift and recalibration requirement.}
The 2026-03-15 follow-up session revealed a 23\% decrease in $\nuZZ$
(3.600\,kHz $\to$ 2.778\,kHz) and a 49\% decrease in $T_2$ (261 $\to$ 133\,$\mu$s)
over nine days. GPS measurements using the stale $\nuZZ$ value produced
feedforward overcorrection of $\sim30\%$ and inflated GPS estimates; correcting
for the measured $\nuZZ$ recovered the expected peak at $\phi^* \approx 0.873$.
This confirms a practical requirement: $\nuZZ$ must be measured via
ZZ Ramsey spectroscopy at the start of each session.

\paragraph{Dephasing-type ZZ in neighbouring qubit pairs.}
Ramsey spectroscopy on Q91--Q92 (ibm\_marrakesh, same session) revealed
$\nuZZ = 33.76\,\mathrm{kHz}$ with a monotone-decreasing correlator $C(\tau)$
and no oscillation---identical in character to the IQM Garnet failures
(Section~\ref{sec:iqm}). This extends the dephasing-type failure mode to a
second device context, demonstrating that dephasing-type ZZ can arise even in
IBM Heavy-Hex processors depending on the specific qubit pair.


% =============================================================
\section{Conclusion}

We have presented EEDT, a ZZ-synchronization-based phase-storage protocol
with analytic design principles and statistically confirmed hardware validation.

\textbf{Key results:}
\begin{enumerate}[leftmargin=*, itemsep=1pt, topsep=2pt]
  \item $\GPS$ factorization derived from Lindblad master equation:
    sync $\times$ decay $\times$ backaction (Eq.~\ref{eq:factorization}).
  \item Synchronization condition: $\taustar = T_2/(\oZZ T_2+1)$
    (Remark~\ref{rem:sync}), with $<5\%$ agreement for ibm\_marrakesh.
  \item EEDT Separation property: $\taustar$ is $N$-independent under
    $\Tmc/\tau_\mathrm{ZZ} < 2\%$ (Analytical Result~\ref{thm:separation}).
  \item Hardware validation: $\GPS = +0.037\pm0.011$ at $\tau=30\,\mu$s,
    combined $z = 8.5\sigma$ (8/10 Holm-corrected conditions).
  \item Multi-device operating predictor: T$_2 > 200\,\mu$s,
    $\eps_{01} < 2\%$, moderate $\nuZZ$ are necessary and sufficient
    conditions (Table~\ref{tab:devices}).
  \item Poisson MLE interpretation: $f_\mathrm{sync}(\phi)\equiv
    P_\mathrm{Poisson}(k{=}1;\lambda{=}\phi)$; $\phi^*=1/A$ is the
    T$_2$-corrected maximum-likelihood optimum
    (Proposition~\ref{prop:mle}).
  \item Lambert $W$ inversion: $\tau_{\pm}(C)=-W_k(-AC/A_s)/(A\oZZ)$
    gives a closed-form design formula and establishes
    $G_\mathrm{max}=A_s/(eA)$ as the analytic protocol bound
    (Remark~\ref{rem:lambertw}).
  \item IQM Garnet comparison: coherent ZZ is a necessary condition for EEDT;
    dephasing-type residual ZZ ($\nuZZ=35$--$46$\,kHz) yields $\GPS<0$
    regardless of feedforward sign, establishing an architectural
    device-selection criterion (Section~\ref{sec:iqm}).
  \item New stochastic synchronisation class: $\phi^*=1$ (Poisson MLE criterion)
    differs from Rabi ($\pi$), Ramsey ($\pi/2$), and CPMG ($\pi$),
    defining stochastic quantum synchronisation as a distinct class.
\end{enumerate}

All code, data, and IBM Quantum job IDs are available at
\href{https://github.com/okudat9}{github.com/okudat9} and
Zenodo~\cite{EEDTv8}.
Contact: \href{mailto:o93dice@gmail.com}{o93dice@gmail.com}.
\emph{License: AGPL-3.0. Commercial use requires prior permission.}


% =============================================================
\section*{Acknowledgments}

The author acknowledges use of IBM Quantum hardware resources and
thanks the open-source Qiskit community.
AI assistance (Claude, Anthropic) was used to accelerate numerical
computation and LaTeX preparation; all physical interpretations and
final manuscript decisions were made by the author.


% =============================================================
\begin{thebibliography}{99}

\bibitem{Preskill2018}
J.~Preskill,
``Quantum computing in the NISQ era and beyond,''
\textit{Quantum} \textbf{2}, 79 (2018).

\bibitem{Arute2019}
F.~Arute et~al.,
``Quantum supremacy using a programmable superconducting processor,''
\textit{Nature} \textbf{574}, 505--510 (2019).

\bibitem{Mundada2019}
P.~S.~Mundada et~al.,
``Suppression of qubit crosstalk in a tunable coupling superconducting circuit,''
\textit{Phys.\ Rev.\ Appl.} \textbf{12}, 054023 (2019).

\bibitem{Krantz2019}
P.~Krantz et~al.,
``A quantum engineer's guide to superconducting qubits,''
\textit{Appl.\ Phys.\ Rev.} \textbf{6}, 021318 (2019).

\bibitem{Temme2017}
K.~Temme, S.~Bravyi, and J.~M.~Gambetta,
``Error mitigation for short-depth quantum circuits,''
\textit{Phys.\ Rev.\ Lett.} \textbf{119}, 180509 (2017).

\bibitem{Endo2018}
S.~Endo, S.~C.~Benjamin, and Y.~Li,
``Practical quantum error mitigation for near-future applications,''
\textit{Phys.\ Rev.\ X} \textbf{8}, 031027 (2018).

\bibitem{Gambetta2012}
J.~M.~Gambetta et~al.,
``Characterization of addressability by simultaneous randomized benchmarking,''
\textit{Phys.\ Rev.\ Lett.} \textbf{109}, 240504 (2012).

\bibitem{Pikovsky1997}
A.~Pikovsky, M.~Rosenblum, and J.~Kurths,
\textit{Synchronization: A Universal Concept in Nonlinear Sciences}
(Cambridge University Press, 2001).

\bibitem{Santos2025}
J.~P.~Santos and R.~Uzdin,
``Drift-resilient mid-circuit measurement error mitigation
for dynamic circuits,''
\textit{arXiv}:2506.11270 (2025).

\bibitem{Koh2024}
J.~M.~Koh et~al.,
``Measurement-based feedforward in superconducting quantum processors,''
\textit{Phys.\ Rev.\ Research} \textbf{6}, 013262 (2024).

\bibitem{Degen2017}
C.~L.~Degen, F.~Reinhard, and P.~Cappellaro,
``Quantum sensing,''
\textit{Rev.\ Mod.\ Phys.} \textbf{89}, 035002 (2017).

\bibitem{Corless1996}
R.~M.~Corless, G.~H.~Gonnet, D.~E.~G.~Hare, D.~J.~Jeffrey, and D.~E.~Knuth,
``On the Lambert $W$ function,''
\textit{Adv.\ Comput.\ Math.} \textbf{5}, 329--359 (1996).

\bibitem{EEDTv8}
T.~Okuda,
``A Self-Stabilizing Qubit Storage Unit with a Defined Operating Window
on IBM Heron r2,''
Zenodo, DOI: \href{https://doi.org/10.5281/zenodo.19029431}%
{10.5281/zenodo.19029431} (2026).

\bibitem{Carmichael1993}
H.~J.~Carmichael,
\textit{An Open Systems Approach to Quantum Optics}
(Springer, Berlin, 1993).

\bibitem{Dalibard1992}
J.~Dalibard, Y.~Castin, and K.~M{\o}lmer,
``Wave-function approach to dissipative processes in quantum optics,''
\textit{Phys.\ Rev.\ Lett.} \textbf{68}, 580--583 (1992).

\bibitem{Breuer2002}
H.~P.~Breuer and F.~Petruccione,
\textit{The Theory of Open Quantum Systems}
(Oxford University Press, Oxford, 2002).

\bibitem{Blais2004}
A.~Blais, R.-S.~Huang, A.~Wallraff, S.~M.~Girvin, and R.~J.~Schoelkopf,
``Cavity quantum electrodynamics for superconducting electrical circuits:
An architecture for quantum computation,''
\textit{Phys.\ Rev.\ A} \textbf{69}, 062320 (2004).

\bibitem{Kuramoto1984}
Y.~Kuramoto,
\textit{Chemical Oscillations, Waves, and Turbulence}
(Springer, Berlin, 1984).

\bibitem{IQMResonance2026}
IQM Resonance cloud quantum computing service,
\url{https://resonance.meetiqm.com}, accessed March 2026.

\bibitem{Krinner2022}
S.~Krinner et~al.,
``Realizing repeated quantum error correction in a distance-three surface code,''
\textit{Nature} \textbf{605}, 669--674 (2022).

\end{thebibliography}


% =============================================================
\appendix

\section{Derivation of \texorpdfstring{$f_\mathrm{sync}$}{f\_sync} from T\texorpdfstring{$_1$}{1} Collapse Statistics}
\label{app:fsync}

Given ancilla collapse rate $\gamma_1 = 1/T_1$, the conditional collapse
time distribution is $p(t^*) = \gamma_1 e^{-\gamma_1 t^*}$
for $t^* \in [0,\tau]$.
The expected ZZ-phase feedforward gain is
\begin{align}
  \mathcal{G}(\tau) &= \int_0^\tau \oZZ\,t^*\,\gamma_1\,e^{-\gamma_1 t^*}\,dt^*
  \notag\\
  &= \frac{\oZZ}{\gamma_1}\bigl[1 - e^{-\gamma_1\tau}(1+\gamma_1\tau)\bigr].
  \label{eq:Gintegral}
\end{align}
For $\gamma_1\tau = \tau/T_1 \ll 1$ (ibm\_marrakesh: $\gamma_1\tau \le 0.123$
for $\tau \le 50\,\mu$s), a Taylor expansion gives
$\mathcal{G}(\tau) \approx \oZZ\tau = \phi$ with
relative error $< \gamma_1\tau/2 < 6.2\%$.
The overcorrection penalty $e^{-\phi}$ arises from phase variance
$\mathrm{Var}[\phi_\mathrm{FF}] \propto \phi$: for $\phi > 1$,
the conditional rotation overshoots with probability $\propto 1 - e^{-\phi}$,
reducing fidelity as $e^{-\phi}$ under a Gaussian approximation valid
to $\mathcal{O}(\phi^2)$.
The product $\mathcal{G}(\tau) \cdot e^{-\phi}
\approx \phi e^{-\phi} = f_\mathrm{sync}$
has its unique maximum at $\phi^* = 1$. $\square$

\section{MCM Dead-Time Kraus Map}
\label{app:deadtime}

During MCM dead-time $\Tmc \approx 700\,\mathrm{ns}$, the target qubit
evolves under amplitude damping and dephasing:
\begin{align}
  K_0 &= \begin{pmatrix}1 & 0 \\ 0 & \sqrt{1-\gamma_1}\end{pmatrix},\quad
  K_1  = \begin{pmatrix}0 & \sqrt{\gamma_1} \\ 0 & 0\end{pmatrix},
\end{align}
with $\gamma_1 = 1-e^{-\Tmc/T_1}$.
For ibm\_marrakesh: $\gamma_1 = 4.1\times10^{-3}$ per MCM shot.
For $N=5$ shots: cumulative dead-time amplitude decay $\approx 2.0\%$,
phase decoherence $\approx 4.6\%$ relative to the $\tau=50\,\mu$s window.
These are the physical origin of $f_\mathrm{back} = e^{-\alpha N\Tmc/\tau}$.

\end{document}
